\section{Conclusion}
\label{sec:conclusion}

SVB's FY2021 annual report disclosed an economic value of equity of \$20,672m and
stated what a 200 basis point rise would do to it: take 27.7\% off. The FY2022
report, filed two weeks before the bank was closed, contains no economic-value figure
at any shock size. This paper asked whether an analyst working only from the earlier
filing could have measured the exposure, and applied the same framework to a bank
that carried comparable securities losses through 2022 and survived.

The answer is that the asset side of the exposure was measurable and the liability
side was not. A model calibrated entirely on 31~December 2021 information, with its
option-adjusted spreads then frozen, priced SVB's held-to-maturity portfolio a year
later at 80.36 per \$100 of carrying amount against the 83.41 the bank disclosed, and
priced Schwab's portfolio to within 1.75\% once its accounting transfers are undone.
The same framework reproduced SVB's published economic-value profile at shock points
it was not fitted to, including the non-monotonicity in its downward shocks, which is
a feature of the shape rather than the level and which a model without a post-shock
floor would miss entirely.

What that exercise surfaces is a property of the disclosure rather than a flaw in it.
An economic-value figure nets an asset valuation the filings constrain tightly against
a deposit valuation they do not constrain at all. Solving for the deposit decay
consistent with SVB's published sensitivity returns an uninsured half-life of
3.59~years. Moving that single unstated parameter to a two-year half-life, no less defensible for
a base that was 96\% uninsured and concentrated in venture-backed technology, more
than doubles the measured exposure, from 27.6\% to
62.5\% of base economic value. Across the full grid of deposit assumptions the figure
runs from $-122\%$ to $+21\%$. The published number was the output of an assumption
about depositor loyalty that the filing did not state and no reader could recover.

The comparison with Schwab locates where that assumption did its work. Both banks
carried similar unrecognised losses into the end of 2022, \$15,152m against
\$14,138m, and both breach the supervisory outlier screen in every bank-year tested.
On the run model only SVB crosses a threshold, at 67.5\% of uninsured deposits, and
it crosses on liquidity rather than solvency. Schwab, funded 18.8\% by uninsured
money against SVB's 95.5\%, crosses neither under any parameter combination tried.
Holding one framework fixed across both, the threshold moves with the uninsured share
rather than with the size of the securities loss. Two observations cannot establish
that funding structure is what separates such banks in general, and they say nothing
about what separated the outcomes in March 2023; \citet{KellyRose2025} make the case
that business model rather than balance sheet is the better discriminator, and
nothing here adjudicates it.

Two things follow for practice rather than for the literature. A disclosed
economic-value sensitivity is interpretable only alongside the deposit assumption that
produced it, and neither bank disclosed one. And the supervisory outlier screen, which
both banks breached in every year tested, separates banks with large rate exposure
from banks without it but says nothing about which of them can survive the
realisation of that exposure. The run threshold and the outlier test answer different
questions, and in 2021 only the second was being reported.
